\section{Method}

\subsection{Overview}

AEGIS-V5 has three main stages:

\begin{enumerate}
\item \textbf{SVD-factorized projections} ( 4.2): each $Q/K/V/O$ projection is
  decomposed with energy-based rank selection and frozen.
\item \textbf{Split-head attention} ( 4.3): heads partitioned into local and global.
\item \textbf{Adaptive gating} ( 4.4): per-head gate fuses local and global.

\end{enumerate}
\begin{lstlisting}
        Input X
           |
            
    +------------------+
    | SVD Q/K/V/O      |  (frozen, energy-ranked)
    | + LoRA adapters  |  (trainable)
    +--------+---------+
             |
    +--------+---------+
                        
 +----------+    +-----------+
 | Local    |    | Global    |
 | Window   |    | RFF       |
 | (2H/3)   |    | (H/3)     |
 +----+-----+    +-----+-----+
      |                |
      +-------+--------+
               
      +---------------+
      | Per-head      |
      | Gating $\sigma$(g)   |
      +-------+-------+
              |
               
         Output
\end{lstlisting}

\subsection{SVD-Factorized Projections}

\subsubsection{Decomposition}

Given a weight $W \in \mathbb{R}^{d_\text{out} \times d_\text{in}}$:

\begin{enumerate}
\item Compute the SVD:
  \[
W = U \Sigma V^\top
\]
\item Select energy-based rank:
  \[
r^* = \min \left\{ r : \frac{\sum\_{i \leq r} \sigma\_i^2}{\sum\_i \sigma\_i^2} \geq \eta \right\}
\]
\item Factor into $W \approx U_r V_r$ where:
  \[
U\_r = U\_{:, :r} \Sigma\_{:r, :r}^{1/2}, \quad
  V\_r = \Sigma\_{:r, :r}^{1/2} V\_{:r, :}^\top
\]

\end{enumerate}
\subsubsection{Parameter Reduction}

\begin{table}[H]
\centering
\small
\begin{tabular}{ll}
\toprule
\textbf{Configuration} & \textbf{Parameters} \\
\midrule
Original & $d_\text{out} \cdot d_\text{in}$ \\
SVD at rank $r$ & $r(d_\text{out} + d_\text{in})$ \\
\bottomrule
\end{tabular}
\end{table}

For $d = 768$ and $r = 384$ (energy 0.95):

\begin{table}[H]
\centering
\small
\begin{tabular}{lll}
\toprule
\textbf{Config} & \textbf{Params} & \textbf{Ratio} \\
\midrule
Original & 589,824 & 1.00$\times$ \\
SVD r=384 & 589,824 & \textbf{1.00$\times$} \\
\bottomrule
\end{tabular}
\end{table}

\textbf{Critical observation.} At $r = 384 = d/2$, SVD parameters equal original. 
For meaningful compression, we need $r < d/2$.

\subsubsection{LoRA on Frozen Factors}

We freeze $(U_r, V_r)$ and add LoRA adapters:

\[
y = U\_r V\_r x + \frac{\alpha}{r\_L} B A x
\]

Only $A, B$ are trained.

\textbf{Note.} AEGIS-V5 uses a \texttt{FactorisedLinear} that applies SVD as a \textbf{stacked} 
composition (\texttt{down} then \texttt{up}), not a reconstruction. This is critical for 
correctness.

\subsection{Split-Head Attention}

\subsubsection{Head Partition}

We split $H$ heads:

\[
H\_\text{local} = \lceil 2H/3 \rceil, \quad
H\_\text{global} = H - H\_\text{local}
\]

For $H = 12$:
\begin{itemize}
\item $H_\text{local} = 8$
\item $H_\text{global} = 4$

\end{itemize}
\subsubsection{Local Branch}

Heads $0 \ldots H_\text{local} - 1$ use sliding window attention:

\[
\text{Out}\_\text{local}[i] = \text{softpick}\!\left(\frac{q\_i K\_\text{local}^\top}{\sqrt{d\_h}}\right) \cdot V\_\text{local}
\]

with a window mask $|i - j| \leq W/2$.

\subsubsection{Global Branch}

Heads $H_\text{local} \ldots H - 1$ use kernel linear attention:

\[
\text{Out}\_\text{global}[i] = \frac{\phi(q\_i)^\top \sum\_j \phi(k\_j) v\_j^\top}{\phi(q\_i)^\top \sum\_j \phi(k\_j)}
\]

with RFF feature map $\phi(x) = \exp(W_\text{rff} x) / \sqrt{F}$.

\subsubsection{Concatenation}

The outputs are concatenated:

\[
\text{Out} = \text{Concat}(\text{Out}\_\text{local}, \text{Out}\_\text{global}) \in \mathbb{R}^{B \times H \times L \times d\_h}
\]

\subsection{Adaptive Gating}

\subsubsection{Per-Head Gate}

Each head has a learnable scalar $g_h$:

\[
\text{Out}\_\text{fused}[i] = \sigma(g\_h + 2.0) \cdot \text{Out}\_\text{local}[i] + (1 - \sigma(g\_h + 2.0)) \cdot \text{Out}\_\text{global}[i]
\]

where $g_h$ is initialized to 0 (so $\sigma(2.0) \approx 0.88$).

\subsubsection{Why Bias 2.0?}

At initialization, we want the model to prioritize local attention (known 
to be safe). As training progresses, the gate can shift toward global.

\subsection{Decoder Fallback}

For decoder models, the local+global split \textbf{does not preserve causality}:

\begin{itemize}
\item Local branch: causal window is fine.
\item Global branch: prefix scan preserves causality.

\end{itemize}
\textbf{But:} For very short sequences ($N < W$), the local branch covers the 
whole sequence, making the split redundant. We fall back to standard dense 
causal attention:

\[
\text{if } \text{is\_decoder} : \text{standard causal softmax}
\]

This preserves correctness at the cost of quadratic complexity for decoders.

\subsection{Full Algorithm}

\textbf{Algorithm 1:} AEGIS-V5 forward
\begin{lstlisting}
Input:  X $\in$ R^{B$\times$L$\times$d}, attention\_mask, is\_decoder
Output: Y $\in$ R^{B$\times$L$\times$d}

1:  Q $\leftarrow$ SVD\_Q(X); K $\leftarrow$ SVD\_K(X); V $\leftarrow$ SVD\_V(X)
2:  if RoPE: Q, K $\leftarrow$ RoPE(Q, K)
3:  if GQA: K, V $\leftarrow$ repeat\_kv(K, V)
4:  if is\_decoder:
5:      # Fallback to standard causal attention
6:      scores $\leftarrow$ Q K^T / $\sqrt{}$d\_h
7:      scores $\leftarrow$ mask\_causal(scores)
8:      out $\leftarrow$ softmax(scores) @ V
9:  else:
10:     # Split-head hybrid attention
11:     Q\_local $\leftarrow$ Q[:, :H\_local]; K\_local $\leftarrow$ K[:, :H\_local]; V\_local $\leftarrow$ V[:, :H\_local]
12:     Q\_global $\leftarrow$ Q[:, H\_local:]; K\_global $\leftarrow$ K[:, H\_local:]; V\_global $\leftarrow$ V[:, H\_local:]
13:     Out\_local $\leftarrow$ local\_window(Q\_local, K\_local, V\_local, W)
14:     Out\_global $\leftarrow$ kernel\_linear(Q\_global, K\_global, V\_global,  )
15:     # Per-head gate
16:     g $\leftarrow$ $\sigma$(blend\_logits + 2.0)
17:     Out $\leftarrow$ g   Out\_local + (1 - g)   Out\_global
18: Y $\leftarrow$ SVD\_O(Out)
19: return Y
\end{lstlisting}

\subsection{Complexity Analysis}

\subsubsection{Encoder}

\begin{table}[H]
\centering
\small
\begin{tabular}{ll}
\toprule
\textbf{Component} & \textbf{Complexity} \\
\midrule
SVD projections & $\mathcal{O}(L r (d + d_h))$ \\
Local attention & $\mathcal{O}(L W d_h \cdot H_\text{local})$ \\
Kernel attention & $\mathcal{O}(L F d_h \cdot H_\text{global})$ \\
Output projection & $\mathcal{O}(L r d)$ \\
\textbf{Total} & $\mathcal{O}(L(W H_\text{local} + F H_\text{global}) d_h)$ \\
\bottomrule
\end{tabular}
\end{table}

For fixed $W, F, H_\text{local}, H_\text{global}$: linear in $L$.

\subsubsection{Decoder}

Quadratic fallback: $\mathcal{O}(L^2 d_h)$.

\subsection{Parameter Analysis}

\begin{table}[H]
\centering
\small
\begin{tabular}{llll}
\toprule
\textbf{Component} & \textbf{Baseline} & \textbf{AEGIS-V5} & \textbf{Change} \\
\midrule
Q/K/V/O (SVD) & $4 d^2$ & $4 r(d + d_h)$ \& depends on $r$ \\
LoRA adapters & 0 & $4 r_L (d + d_h)$ & +small \\
RFF proj & 0 & $H_\text{global} \cdot d_h \cdot F$ & +medium \\
blend\_logits & 0 & $H$ & +negligible \\
\bottomrule
\end{tabular}
\end{table}

\textbf{Critical.} Whether AEGIS-V5 reduces parameters depends entirely on $r$.

For $d = 768$, $r = 384$ (energy 0.95 default):
\begin{itemize}
\item SVD: no compression (ratio 1.00$\times$)
\item LoRA: +98,304
\item RFF: $4 \cdot 64 \cdot 128 = 32,768$
\item \textbf{Net: +131,072 params per layer}

\end{itemize}
For $d = 768$, $r = 128$:
\begin{itemize}
\item SVD: $128 \cdot 1536 = 196,608$ (vs 589,824) $\to$ \textbf{3$\times$ compression}
\item LoRA: +98,304
\item RFF: +32,768
\item \textbf{Net: -262,144 params per layer}

\end{itemize}
\textbf{We are explicit:} The parameter effect depends on the SVD rank.

\subsection{Theoretical Notes}

\subsubsection{SVD Approximation Error}

If $W = U\Sigma V^\top$, then:

\[
\|W - U\_r \Sigma\_r V\_r^\top\|\_F = \sqrt{\sum\_{i \textgreater{} r} \sigma\_i^2}
\]

The energy-based rank ensures this is below $\sqrt{1 - \eta} \|W\|_F$.

\subsubsection{Kernel Approximation}

Standard RFF bound (Rahimi \& Recht, 2007):

\[
\sup\_{x, y} |\phi(x)^\top \phi(y) - e^{x \cdot y}| \leq \epsilon
\]

with $F = \mathcal{O}(d_h \log(1/\epsilon^2))$.

\subsubsection{Causality}

\textbf{Claim.} For decoders, the fallback preserves causality.

\textbf{Proof.} Fallback uses standard causal softmax with triangular mask. $\square$

\textbf{Claim.} For encoders, the split-head design preserves non-causality 
(no masking needed).

\textbf{Proof.} Local window covers $|i - j| \leq W/2$; kernel attention covers 
all $j$. Both are bidirectional by design. $\square$
\begin{lstlisting}

---
\end{lstlisting}